% ------------------------------------------------------------------
%  Capitolo 12 — Immagini, parole e palazzi della memoria
%  Parte III
%  Ricerca di riferimento: ricerca 01 §2.4, §5.6–5.7, §6.1, §14.1;
%    03 §3.2–3.3, §4.4
%  Voci bibliografiche principali: paivio1986, mayer2020, wammes2016,
%    dresler2017, spence1984, dunlosky2013
%  Epigrafe: ✔ verificata sul testo (K-17 chiuso: III, 16, 29)
%  Scheda del capitolo: ricerca/06_indice-ragionato.md, capitolo 12
%
%  STATO: prima stesura (23/09/2026), circa 2.700 parole.
%  Dati verificati: S-13 (Dresler et al. 2017: 40 sessioni quotidiane di
%  30 minuti in sei settimane; da circa 26–30 a 62 parole su 72; controllo
%  attivo +7; effetti ancora presenti a quattro mesi). Wammes et al. 2016
%  verificato il 23/09/2026 sul comunicato dell'Università di Waterloo
%  (40 secondi per parola; spesso più del doppio delle parole ricordate;
%  qualità del disegno irrilevante; vantaggio anche con 4 secondi).
%  Rhetorica ad Herennium III, 22, 35 ✔ (parafrasato); III, 22, 37 non
%  citato in latino per via della variante mutas/multas.
% ------------------------------------------------------------------
\chapter{Immagini, parole e palazzi della memoria}
\label{cap:immagini}

\epigrafe[La memoria artificiale consiste dunque di luoghi e di immagini.]{Constat igitur artificiosa memoria ex locis et imaginibus.}{Rhetorica ad Herennium III, 16, 29}

\iniziale{V}{erso l'anno} 85 prima di Cristo, un maestro di retorica di
cui non conosciamo il nome scrive per un certo Gaio Erennio un manuale
dell'arte del parlare. Nel terzo libro, dopo aver spiegato come si
costruisce e si pronuncia un discorso, affronta un problema pratico:
come fa un oratore a ricordare un'orazione lunga ore, senza appunti?
La risposta comincia con un'osservazione di buon senso. Le cose che
vediamo ogni giorno, piccole, abituali, consuete, di solito non le
ricordiamo, perché la mente non trova in esse nulla che la colpisca;
ricordiamo invece a lungo ciò che è eccezionale, vergognoso, inaudito,
grandioso, incredibile o ridicolo.\footnote{\emph{Rhetorica ad
Herennium} III, 22, 35.} Chi ricorda la colazione di martedì di tre
settimane fa? Eppure tutti ricordiamo il giorno in cui qualcosa, a
colazione, è andato storto in modo memorabile.

Da questa osservazione l'anonimo autore ricava un metodo:
\enquote{la memoria artificiale consiste di luoghi e di immagini}.
Bisogna immaginare una serie ordinata di luoghi ben noti e collocare in
ciascuno un'immagine vivida, strana, in movimento, che rappresenti ciò
che si vuole ricordare. Poi, al momento di parlare, basta percorrere
mentalmente i luoghi, uno dopo l'altro, e ritrovare le immagini lasciate
lungo la strada.

Duemila anni dopo, questo metodo è ancora quello usato dai campioni delle
gare di memoria, capaci di memorizzare in pochi minuti
l'ordine di un mazzo di carte o centinaia di cifre. Nel 2017 un gruppo
di neuroscienziati guidato da Martin Dresler ha voluto capire se questi
campioni avessero un cervello diverso dagli altri, o soltanto un
allenamento diverso \parencite{dresler2017}. Ha preso persone comuni,
senza alcuna esperienza di mnemotecniche, e ne ha allenate alcune con
il metodo dei luoghi: mezz'ora al giorno, per quaranta sessioni, nell'arco
di sei settimane. Prima dell'allenamento, da una lista di settantadue
parole ne ricordavano in media meno di trenta. Dopo, ne ricordavano
sessantadue, mentre un gruppo che si era allenato con un altro tipo di
esercizi di memoria era migliorato di sole sette parole. Quattro mesi
più tardi, senza altro allenamento, il vantaggio c'era ancora; e le
immagini del loro cervello cominciavano a somigliare a quelle dei
campioni.

È un risultato spettacolare, e giustamente celebre. Ma questo capitolo
deve rispondere a una domanda più modesta e più utile: le immagini, i
disegni e le tecniche di memoria servono \emph{a chi studia}? La
risposta è sì, ma con alcune distinzioni importanti.

\section{Parole e immagini}

Negli anni Settanta lo psicologo canadese Allan Paivio propose una
teoria semplice: la mente elabora le informazioni attraverso due
sistemi distinti ma collegati, uno per il linguaggio e uno per le
immagini \parencite{paivio1986}. Una parola astratta, come
\emph{giustizia}, si affida quasi soltanto al sistema verbale; una
parola concreta, come \emph{bilancia}, attiva anche un'immagine. Le
informazioni codificate in entrambe le forme hanno due vie per essere
ritrovate, e per questo si ricordano meglio. È la \emph{teoria della
doppia codifica}, e spiega un fatto che tutti conosciamo: le parole
concrete e le immagini si ricordano più facilmente delle parole
astratte.

Da allora Richard Mayer, dell'Università della California a Santa
Barbara, ha condotto e raccolto centinaia di esperimenti su come si
impara da testi e immagini insieme, e ne ha ricavato una serie di
principi pratici \parencite{mayer2020}. Il primo è il più generale: si
impara meglio da parole e immagini che da sole parole. Gli altri dicono
\emph{quali} immagini e \emph{come}. Le immagini devono essere
pertinenti: un diagramma che mostra come funziona una pompa aiuta a
capire la pompa; una fotografia decorativa di un'officina no, e anzi
distrae. Testo e immagine devono stare vicini, sulla stessa pagina,
con le etichette sul disegno, non in una legenda lontana che obbliga
l'occhio a fare avanti e indietro. E le decorazioni, i dettagli
suggestivi ma irrilevanti, vanno eliminati: tutto ciò che non serve a
capire consuma una parte dell'attenzione che serve a capire.

Per chi studia, questi principi hanno due conseguenze. La prima
riguarda i materiali: quando un argomento ha una struttura spaziale o
temporale -- un'anatomia, un circuito, un processo chimico, una
sequenza di eventi, un'organizzazione istituzionale -- cerca o costruisci
una rappresentazione visiva, e studia testo e immagine insieme. La
seconda riguarda gli appunti: molte informazioni si capiscono e si
ricordano meglio se le trasformi in una forma visiva adatta al loro
contenuto: una linea del tempo per la storia, un diagramma di flusso per
una procedura, una tabella di confronto per istituti o malattie simili,
uno schema ad albero per una classificazione. Non si tratta di rendere
gli appunti belli, ma di renderli \emph{strutturati}: la forma visiva
deve mostrare le relazioni tra le cose.

Lo stesso vale per le mappe concettuali e le mappe mentali, così
diffuse nella scuola italiana. Come strumento per organizzare un
argomento sono utili; ma nel capitolo~\ref{cap:recupero} abbiamo visto
che, a parità di tempo, ricordare un testo a libro chiuso fa imparare
di più che costruirne la mappa con il testo davanti. La soluzione non è
rinunciare alle mappe: è costruirle \emph{a memoria}. Disegna la mappa
di un capitolo senza guardarlo, poi apri il libro, confronta e
completa con un altro colore. La parte aggiunta in colore è ciò che
devi ristudiare. Una mappa fatta così non è più un prodotto da
decorare, ma un esercizio di recupero con una forma visiva.

Un'ultima abitudine da prendere riguarda le figure dei manuali. Molti
studenti le saltano, o le guardano di sfuggita, perché \enquote{il testo
dice già tutto}. Nelle discipline scientifiche e tecniche è spesso
vero il contrario: la figura dice in un colpo d'occhio ciò che il testo
impiega una pagina a spiegare, e gli autori l'hanno costruita con cura
proprio per questo. Fermati sulle figure, leggi la didascalia, spiegati
che cosa rappresenta ogni parte e come si collega al testo: è
l'auto-spiegazione del capitolo~\ref{cap:capire}, applicata a
un'immagine.

Tommaso d'Aquino, nelle sue quattro regole per ricordare, lo aveva
detto a modo suo. La prima regola chiede di associare a ciò che si
vuole ricordare \enquote{immagini adatte, ma non del tutto consuete,
perché ciò che è insolito ci meraviglia di più, e così l'animo vi si
trattiene di più e con più forza}; la seconda, di disporre
\enquote{ordinatamente} ciò che si vuole tenere a
memoria.\footnote{Tommaso d'Aquino, \emph{Summa theologiae} II-II,
q.~49, a.~1, ad~2. Le altre due regole, sull'attenzione e sulla
meditazione frequente, le abbiamo incontrate nei
capitoli~\ref{cap:recupero} e~\ref{cap:tempo}.} Immagini e ordine:
proprio ciò che la psicologia avrebbe misurato sette secoli più tardi.

Un'avvertenza, però, va fatta subito. La doppia codifica \emph{non}
significa che esistano \enquote{persone visive} che imparano meglio con
le immagini e \enquote{persone uditive} che imparano meglio ascoltando.
È una delle convinzioni più diffuse nella scuola e nell'università, e
non ha il sostegno delle prove: adattare l'insegnamento allo
\enquote{stile} dichiarato da ciascuno non migliora l'apprendimento. Ne
parleremo nel capitolo~\ref{cap:miti}. Ciò che conta non è lo stile
della persona, ma la natura del contenuto: la geografia si impara con
le carte geografiche, per tutti; la musica con l'ascolto, per tutti.
Le buone immagini aiutano tutti, quando il contenuto si presta.

\section{Disegnare per ricordare}

Se guardare un'immagine aiuta, farla aiuta ancora di più. Nel 2016
Jeffrey Wammes, Melissa Meade e Myra Fernandes, dell'Università di
Waterloo, in Canada, hanno dato ad alcuni studenti una serie di parole
semplici e concrete, come \emph{mela}, con un compito diverso per
ciascuna: per alcune, disegnare ciò che la parola indicava, in quaranta
secondi; per altre, scriverla ripetutamente per lo stesso tempo
\parencite{wammes2016}. Dopo un compito di distrazione, dovevano
ricordare quante più parole possibile. Le parole disegnate venivano
ricordate molto meglio, spesso più del doppio di quelle scritte. Il
vantaggio si manteneva anche con soli quattro secondi per disegno, e
batteva anche altre strategie, come elencare le caratteristiche
dell'oggetto, immaginarlo o guardarne una fotografia. E soprattutto: la
qualità del disegno non contava. Uno scarabocchio funzionava come
un'opera d'arte.

Perché? Disegnare costringe a fare molte cose insieme: pensare al
significato della parola, decidere come rappresentarlo, muovere la mano,
guardare il risultato. Il ricordo che ne nasce è insieme verbale,
visivo e motorio, e porta il segno di un piccolo sforzo di elaborazione.
È, in un certo senso, un effetto di generazione
(capitolo~\ref{cap:capire}) in forma grafica.

Anche qui, però, serve onestà. Gli esperimenti di Wammes riguardano
parole semplici e facili da disegnare; le prove che il disegno aiuti
allo stesso modo con contenuti complessi e astratti sono ancora poche.
Nessuno sa bene come disegnare \enquote{il principio di sussidiarietà}.
Il consiglio pratico, quindi, è mirato: disegna ciò che ha una forma o
una struttura -- un organo, una cellula, un meccanismo, un apparato, un
grafico, lo schema di un procedimento -- e disegnalo, quando puoi,
\emph{a memoria}. Chiudi il libro e ridisegna il ciclo di Krebs, la
struttura del nefrone, lo schema del processo civile; poi confronta con
l'originale e correggi. Un disegno fatto a memoria è pratica del
recupero, doppia codifica e generazione, tutte insieme.

\section{Il metodo dei luoghi}

Torniamo ai luoghi e alle immagini dell'anonimo autore. Il metodo ha
una storia leggendaria, che abbiamo raccontato nel
capitolo~\ref{cap:storia}: il poeta Simonide che, sopravvissuto al
crollo della sala di un banchetto, riesce a identificare le vittime
ricordando il posto in cui ciascuna era seduta, e ne ricava che è
soprattutto l'ordine a dare luce alla memoria. Per secoli, da Cicerone
e Quintiliano ai predicatori medievali e ai maestri del Rinascimento,
l'\emph{ars memoriae} è stata una parte essenziale della formazione di
chiunque dovesse parlare in pubblico.

Uno degli episodi più singolari di questa storia ha per protagonista
un italiano. Matteo Ricci, gesuita di Macerata, arrivato in Cina alla
fine del Cinquecento, stupiva i letterati cinesi con la sua memoria e
insegnava loro l'arte dei \enquote{palazzi della memoria}, che
presentava come un aiuto per preparare i durissimi esami imperiali, sui
quali si giocava l'accesso alla carriera di funzionario. La vicenda è
raccontata dallo storico Jonathan Spence in un libro che ha reso celebre
l'espressione \emph{palazzo della memoria} \parencite{spence1984}.

Come si costruisce un palazzo della memoria? Si procede in tre passi.

\emph{Scegli i luoghi.} Pensa a un percorso che conosci benissimo: la
tua casa, la strada che fai ogni giorno per andare all'università, il
tuo paese. Individua lungo il percorso una serie di punti precisi, in
un ordine fisso: la porta d'ingresso, l'attaccapanni, il divano, la
finestra del soggiorno, il frigorifero, e così via. I luoghi devono
essere distinti tra loro e sempre percorsi nello stesso ordine.

\emph{Crea le immagini.} Per ogni elemento da ricordare, inventa
un'immagine concreta, esagerata, possibilmente in movimento: come diceva
già la \emph{Rhetorica ad Herennium}, le immagini che fanno qualcosa, e
che sono insolite, si ricordano meglio di quelle immobili e banali. Più
è strana, meglio è.

\emph{Colloca e ripercorri.} Metti ogni immagine in un luogo, nell'ordine
della lista. Poi percorri mentalmente il palazzo, più volte e a
distanza di tempo, e ritrova le immagini. Qui il metodo antico incontra
quelli di questo libro: il palazzo non si costruisce una volta per
tutte, si \emph{ripercorre}, cioè si recupera, e a intervalli.

Il metodo dei luoghi è straordinariamente efficace per un tipo preciso
di materiale: le \emph{sequenze ordinate} di elementi che non hanno un
legame logico tra loro. Serve per ricordare l'ordine degli argomenti di
un discorso o di una presentazione, la scaletta di una risposta
all'esame orale, le fasi di un procedimento, una serie di nomi. Per
questi usi è uno strumento potente, che chiunque può imparare in poche
ore di pratica, come ha mostrato l'esperimento di Dresler.

\begin{esercizio}
\textbf{Costruisci un palazzo della memoria.} Scegli un elenco ordinato
che devi imparare: le fasi della mitosi, i re di Roma, i passaggi di un
procedimento amministrativo, gli argomenti di una relazione che dovrai
presentare. Per le fasi della mitosi, per esempio: sulla porta di casa
un \emph{professore} in toga che fa lezione (profase); sull'attaccapanni
lo stesso professore che \emph{promette} a gran voce qualcosa agli
studenti (prometafase); sul divano, esattamente a \emph{metà}, una fila
di ospiti seduti in riga (metafase: i cromosomi allineati al centro);
alla finestra un \emph{ananas} tagliato in due, le cui metà volano
verso i due lati opposti della stanza (anafase: i cromatidi che si
separano); davanti al \emph{televisore} due famiglie identiche, ciascuna
con il proprio divano (telofase: due nuclei). Percorri il palazzo tre volte
oggi, una domani, una tra una settimana. Poi prova a ricordare
l'elenco senza percorrerlo.
\end{esercizio}

\section{Mnemotecniche: quando servono e quando no}

Il metodo dei luoghi è la più nobile delle \emph{mnemotecniche}, ma non
la sola. Ci sono gli acronimi e le frasi in cui le iniziali delle parole
richiamano un elenco; ci sono le rime e le filastrocche; e c'è il
metodo della \emph{parola chiave}, utile soprattutto per i vocaboli
stranieri: per ricordare che il tedesco \emph{Tisch} significa
\enquote{tavolo}, si può immaginare una tisana bollente rovesciata su un
tavolo apparecchiato. Tutte queste tecniche funzionano
allo stesso modo: creano un aggancio artificiale per un'informazione
che, di per sé, non ne ha.

Ed è qui che si vede il loro limite. Le mnemotecniche aiutano a
\emph{ritrovare} un'informazione, non a \emph{capirla}. Sanno dirti
qual è la terza fase della mitosi, non che cosa succede in quella fase
né perché. Per questo la rassegna di Dunlosky e colleghi, che valuta le
tecniche di studio per la loro utilità generale, assegna alla parola
chiave un'utilità bassa: funziona bene per un tipo ristretto di
materiale, soprattutto vocaboli, e i suoi effetti possono svanire presto
se non si ripassa \parencite{dunlosky2013}. Chi volesse studiare un
manuale di diritto costituzionale o di biochimica con le mnemotecniche
costruirebbe un magnifico palazzo pieno di parole senza significato.

La regola pratica è dunque questa. Usa le mnemotecniche per ciò che è
davvero \emph{arbitrario}: elenchi, sequenze, nomenclature, vocaboli,
associazioni senza logica interna. Per tutto il resto -- cioè per la
maggior parte di ciò che si studia all'università -- la strada passa per
la comprensione, i collegamenti, il recupero e la distribuzione dei
capitoli precedenti. E anche le mnemotecniche, per durare, hanno bisogno
di essere riprese nel tempo: un palazzo della memoria non visitato si
svuota come qualunque altro ricordo.

Diffida, infine, di chi promette di \enquote{memorizzare un libro in un
giorno} o di \enquote{triplicare la memoria in una settimana}. I
campioni di memoria esistono, e il metodo dei luoghi funziona; ma i
campioni si allenano per mesi e anni, e ciò che ricordano sono per lo
più liste di elementi arbitrari, non conoscenze da capire e da usare.
Lo stesso autore anonimo della \emph{Rhetorica ad Herennium}, alla fine
della sua lunga esposizione, avvertiva che in ogni disciplina
\enquote{l'insegnamento della tecnica è debole senza la massima
assiduità dell'esercizio}.\footnote{\emph{Rhetorica ad Herennium} III,
24, 40. È l'epigrafe del capitolo~\ref{cap:miti}.} Nessun trucco
dispensa dalla fatica di studiare.

\asterismo

Le immagini, dunque, aiutano, e aiutano tutti: quando sono pertinenti,
quando stanno vicino alle parole, quando le costruisci tu, e ancora di
più quando le ridisegni a memoria. Le tecniche di memoria sono
strumenti nobili e antichi, utilissimi per ciò che è arbitrario e
ordinato. Ma i luoghi e le immagini servono a ritrovare; capire è un
altro lavoro, e ricordare a lungo un altro ancora.

\begin{daricordare}
Parole e immagini insieme si ricordano meglio, e disegnare aiuta a
ricordare. Le immagini aiutano tutti quando il contenuto è visivo: non
esistono \enquote{persone visive}. Le mnemotecniche servono per liste e
associazioni arbitrarie, non per capire.
\end{daricordare}

\begin{domande}
  \item Che cosa dice la teoria della doppia codifica? Perché non
        conferma l'esistenza di \enquote{persone visive}?
  \item Quali caratteristiche deve avere un'immagine per aiutare a
        capire un testo, secondo i principi di Mayer?
  \item Che cosa hanno mostrato Dresler e colleghi allenando persone
        comuni con il metodo dei luoghi?
  \item Perché le mnemotecniche non bastano per studiare? Per quali
        contenuti di una materia che studi potrebbero servire?
  \item (Ripresa, capitolo~\ref{cap:capire}) Che cos'è
        l'auto-spiegazione, e perché funziona?
  \item (Ripresa, capitolo~\ref{cap:storia}) Che cosa ricavò Simonide
        dalla tragedia del banchetto?
\end{domande}

\begin{sintesi}
  \item Le informazioni codificate sia in parole sia in immagini hanno
        due vie di recupero e si ricordano meglio.
  \item Le immagini utili sono pertinenti, vicine al testo e prive di
        decorazioni; negli appunti, la forma visiva deve mostrare le
        relazioni tra le cose.
  \item Disegnare ciò che si impara, anche male, migliora il ricordo;
        disegnare a memoria unisce disegno e recupero.
  \item Il metodo dei luoghi è efficacissimo per le sequenze ordinate e
        si impara in poche settimane; le altre mnemotecniche servono per
        vocaboli e associazioni arbitrarie.
  \item Le mnemotecniche aiutano a ritrovare, non a capire: sono
        strumenti di nicchia, non un metodo di studio.
\end{sintesi}

\finecapitolo
